\section{Formative Study}
To gain insight into how users use current natural language-based frameworks to coordinate multiple LLM-based agents and identify the challenges that exist during the exploration process for coordination strategy, we carried out a formative study. Based on the findings in the formative study, we formulated four design requirements to enhance the process of designing coordination strategies for LLM-based multi-agent collaboration.
\subsection{Participants and Procedure}

We recruited 8 participants who have experience or general interest in LLM-based multi-agent collaboration from the local university and online discussion platforms for open-source multi-agent frameworks. Four of them are experienced experts in LLM-based multi-agent systems (E1 and E2 are NLP researchers familiar with LLM-based multi-agent collaboration, while E3 and E4 are developers having experience in constructing multi-agent systems). Another four of them (G1-4) are general users who have a basic understanding of LLM-based agents and are interested in building their own LLM-based multi-agent collaboration strategy. 

\textbf{Procedure:} In our formative interviews,
we initially asked the participants about their prior experience with any LLM-based multi-agent framework or system. Afterward, we show participants how to use natural language to specify the coordination strategy for multi-agent collaboration using AutoAgents \cite{AutoAgents} alongside a text editor and the ``group chat mode'' of AutoGen\cite{wu2023autogen}. After the participants got familiar with the usage, they were asked to choose a task that could benefit from collaboration among multiple LLM-based agents and construct their own coordination strategy with both systems separately. The participants can also use ChatGPT to assist them in designing coordination strategies during the process. Finally, we gathered feedback on participants' experience during the construction of the coordination strategy and inquiry about the challenges they faced during the process. For participants who reported prior experience with any LLM-based multi-agent framework at the start of the interview, we also asked them to compare the strengths and weaknesses of natural language-based methods with previous frameworks they have used.

% on the strengths and weaknesses of the natural language based  they use these frameworks
\subsection{Findings}
All participants find using natural language to design coordination strategies is intuitive and could be a promising approach to democratizing agent coordination for a wider general audience. However, several challenges are also identified, which hinder the participants' current exploration process to design coordination strategies at their will.

\textbf{Lack of structure to regularize the ambiguity of natural language.} 
Although natural language is easy to understand and highly expressive, it is prone to ambiguity. For example, for a high-level cooperation strategy description ``\textit{Pharmaceutical Chemist, Patent Agent, and Clinical Research Scientist collaboratively drafting a patent application}'', there could be many ambiguous aspects: ``\textit{Who takes the main responsibility for drafting?}'', ``\textit{How are opinions integrated?}''... During the formative study, we notice that users often start with a generated collaborative strategy and then, upon observing the unexpected outcomes, identify areas that were not articulated and make remedial enhancements to the original cooperation strategy. After several rounds, the collaborative strategy specification ``\textit{could be lengthy and messy to read}'' (G4), and ``\textit{sometimes even contain self-conflicts}'' (G3). Both E1 and E4 suggest that ``\textit{some structures should be provided to regularize the design process}'' (E1, E4), which can ``\textit{draw on designs from current code-based frameworks}'' (E1).

\textbf{Lost in the vast amount of intricate text.} During the process of designing collaborative strategies, users need to refer to a substantial amount of text information (e.g., previously designed collaborative strategies, descriptions of different agents, input/output of agents, intermediate objects). The vast amount of intricate text poses significant cognitive overhead on users during the design process. Many participants express the feeling that the quantity of text is ``\textit{overwhelming}'' (E3, E4, G1-G4). They usually needed to ``manually switch back and forth between different parts of the texts'' (G2) and ``sometimes forget where to find'' (G3). E3 mentioned that ``\textit{as the complexity of my strategy rises, maintain a clear connection between specific execution result and the corresponding part of my strategy become challenging}'' (E3). The Participants also express the need to ``\textit{have a visual interface that helps organize information}'' (G2, E3).

\textbf{Lack of interactions support to facilitate exploration.} To leverage the powerful coordination capabilities of LLMs and deal with ``writer’s block'', participants often chat with LLMs (e.g. using ChatGPT) to aid them in drafting and exploring coordination strategies. However, the linear non-reversible conversation interface for chatting is not designed for iterative multi-thread exploration. E2 mentioned that  ``\textit{managing exploration history and toggling between different possibilities with manual copy \& paste is cumbersome}'' (E2). Additionally, participants also mentioned concerns over the fluidity of exploration due to the need for ``\textit{manually craft auxiliary prompts for different exploration purposes}'' (G2). Moreover, due to the stochastic nature of LLM outputs (which is also greatly affected by prompts) and diverse possibilities for strategy design, participants express the need for ``\textit{an interface to help systematically explore and compare different outputs by LLM}'' (G1). 

\subsection{Design Requirement}
In response to the problems identified in the formative study, our goal is to develop an interactive system to help general users smoothly explore and design coordination strategies for LLM-based Multi-agent collaboration. The design requirements are summarized as follows:

\textbf{R1: Generate a structured coordination strategy for the user's goal.} Although using natural language to describe a coordination strategy lowers the barrier for users and brings a high level of flexibility, users often lack an effective way to deal with its ambiguity. Therefore, the system should provide a structure of coordination strategy to help regularize the ambiguity inherent in natural language and serve as a scaffolding for downstream exploration. Moreover, to help the user kick off, the system should be able to generate an initial coordination strategy based on the user's goal leveraging the coordination capability of LLMs.

\textbf{R2: Provide an effective visual organization for the strategy.} When users are devising a coordination strategy, they need to refer to various types of relevant information represented in the text. However, at present, users have to flip back and forth through a vast amount of plain text to search for target information and do verification, which creates a significant cognitive load. Therefore, the system should effectively visually organize and enhance the various pieces of information involved in the coordination strategy design process to help users quickly locate the information they need and provide the relevant context.

\textbf{R3: Support flexible interactions to facilitate strategy exploration.} Users often need to explore various possible options at different stages of the coordination strategy design with the help of LLMs; however, iterative exploration based on a linear chat interface with LLMs is cumbersome and unintuitive. Therefore, the system should support flexible and intuitive interactions to help users conduct multi-thread iterative exploration. Moreover, the system needs to offer assistance to help systematically explore and compare different design choices for coordination strategies.

% interactive assistance for users to efficiently invoke and explore the outputs of LLMs.
% \textbf{R4: Establish connections between the strategy and its result.}
\textbf{R4: Provide visual enhancement for the execution result.} During the execution of collaborative tasks, agents generate a substantial amount of textual information. However, at present, both code-based and natural language-based agent coordination frameworks only output results through a plain text terminal. Users often have to manually switch back and forth between different parts of the coordination strategy and the execution results to establish connections, which increases cognitive load and decreases analytical efficiency. Therefore, the system needs to provide visual enhancements to help users examine execution results. 
